\makeatletter
\newcommand{\appendsection}[1]{\vspace{2.75ex plus .15ex minus .05ex}\noindent{\addcontentsline{toc}{subsection}{#1}\large\textbf{#1}}\vspace{.5ex plus .05ex minus .05ex}\nopagebreak[4]\par\noindent\trim@spaces}
\makeatother

\markboth{}{}
\renewcommand{\thepage}{A-\arabic{page}}
\setcounter{page}{1}
%\thispagestyle{plain}
%\noindent
\chapter*{Lampiran}
\addcontentsline{toc}{chapter}{Lampiran}
Matematika tersusun dari argumen (yakni penalaran,
bukan adu lempar pecah-belah).
Bagian ini menguraikan secara ringkas materi latar belakang dan
teknik berargumen yang kita gunakan dalam buku ini.

Bagian ini memberikan garis besar informal tanpa menyertakan pembuktian.
Untuk pembahasan lebih lanjut, \cite{HowToProveIt} sangat baik.
Dua sumber lain yang tersedia daring adalah
\cite{Hefferon} dan~\cite{Beck}.




%\vskip .75in
\appendsection{Pernyataan}
%\addcontentsline{toc}{subsection}{Propositions}
%\bigskip
%\par\noindent
Pernyataan matematika formal diberi label
\definend{Teorema} % \index{theorem} 
untuk hasil utama,
\definend{Korolari} % \index{corollary}  
untuk hasil yang langsung mengikuti hasil sebelumnya, atau
\definend{Lema} %\index{lemma} 
untuk hasil yang terutama digunakan dalam membuktikan hasil lain.

Pernyataan dapat bersifat kompleks dan terdiri atas banyak bagian.
Benar atau salahnya seluruh pernyataan bergantung pada nilai kebenaran
bagian-bagiannya sekaligus pada cara bagian-bagian itu dirangkai.

\startword{Tidak}
Jika \( P \) adalah proposisi,
`tidak benar bahwa \( P \)' bernilai benar apabila \( P \) bernilai
salah.
Sebagai contoh, `\( 6 \) bukan bilangan prima' benar karena \( 6 \)
merupakan hasil kali bilangan-bilangan bulat positif yang lebih kecil.

Untuk membuktikan pernyataan `tidak \( P \)', tunjukkan bahwa \( P \) salah.



\startword{Dan}
Agar pernyataan berbentuk `\( P \) dan \( Q \)' benar,
kedua bagiannya harus benar:
`\( 7 \) bilangan prima dan \( 3 \) juga bilangan prima' benar, sedangkan
`\( 7 \) bilangan prima dan \( 3 \) bukan bilangan prima' salah.

Untuk membuktikan `\( P \) dan \( Q \)', buktikan masing-masing bagiannya.




\startword{Atau}
Pernyataan `\( P \) atau \( Q \)' benar apabila sedikitnya salah satu
bagiannya benar:
`\( 7 \) bilangan prima atau \( 4 \) bilangan prima' benar, sedangkan
`\( 8 \) bilangan prima atau \( 4 \) bilangan prima' salah.
Jika kedua klausa pernyataan itu benar,
seperti dalam `\( 7 \) bilangan prima atau \( 3 \) bilangan prima',
kita menganggap pernyataan secara keseluruhan tetap benar.
(Dalam percakapan sehari-hari orang kadang-kadang memakai `atau' secara
eksklusif\Dash ``Hidup merdeka atau mati'' tidak dimaksudkan agar kedua
bagiannya terjadi\Dash tetapi kita tidak akan memakai `atau' dengan cara itu.)

Untuk membuktikan `\( P \) atau \( Q \)', tunjukkan bahwa dalam setiap kasus
sekurang-kurangnya satu bagian benar (mungkin pada suatu kasus bagian yang
satu dan pada kasus lain bagian yang lain, tetapi selalu setidaknya satu).



\startword{Jika--maka}
\index{pernyataan jika--maka}
Pernyataan `jika \( P \), maka \( Q \)' juga dapat ditulis sebagai
`\( P \) mengimplikasikan \( Q \)', `\( P\implies Q\)',
`\( P \) cukup untuk menghasilkan \( Q \)', atau `$Q$ jika $P$'.
Pernyataan itu benar kecuali ketika \( P \) benar sedangkan \( Q \) salah.
Jadi `jika \( 7 \) bilangan prima, maka \( 4 \) bukan bilangan prima' benar,
sedangkan `jika \( 7 \) bilangan prima, maka \( 4 \) juga bilangan prima' salah.
(Berbeda dari pemakaiannya dalam percakapan sehari-hari, dalam matematika
`jika \( P \), maka \( Q \)' tidak menyiratkan bahwa \( P \) mendahului
atau menyebabkan \( Q \).)

Perhatikan akibat dari paragraf sebelumnya: jika $P$ salah, maka
`jika \( P \), maka \( Q \)' benar tanpa memandang nilai kebenaran $Q$:
`jika \( 4 \) bilangan prima, maka \( 7 \) bilangan prima' dan
`jika \( 4 \) bilangan prima, maka \( 7 \) bukan bilangan prima'
keduanya merupakan pernyataan benar.
(Keduanya \definend{benar secara hampa}.\index{benar secara hampa})
Perhatikan pula bahwa `jika \( P \), maka \( Q \)' benar ketika $Q$ benar:
`jika $4$ bilangan prima, maka $7$ bilangan prima' dan
`jika $4$ bukan bilangan prima, maka $7$ bilangan prima' keduanya benar.

% (We adopt this definition of implication
% because we want statements 
% such as `if $n$ is a perfect square then $n$ is not prime'
% to be true no matter which 
% number $n$ appears in that statement.
% For instance, we want `if $5$ is a perfect square then $5$ is not prime'
% to be true so we want that if both $P$ and~$Q$ are false then
% $P\implies Q$ is true.)

Ada dua cara utama untuk membuktikan suatu implikasi.
Cara pertama bersifat langsung:~andaikan \( P \) benar dan gunakan
asumsi itu untuk membuktikan \( Q \).
Sebagai contoh, untuk menunjukkan `jika suatu bilangan habis dibagi 5,
maka dua kali bilangan itu habis dibagi 10', kita dapat mengandaikan
bilangan tersebut adalah \( 5n \), lalu menyimpulkan \( 2(5n)=10n \).
Cara tidak langsung adalah membuktikan pernyataan
\definend{kontraposisi}\index{kontraposisi}:
`jika \( Q \) salah, maka \( P \) salah'
(dengan kata lain, `\( Q \) hanya dapat salah ketika \( P \) juga salah').
Jadi untuk menunjukkan `jika suatu bilangan asli prima, maka bilangan itu
bukan kuadrat sempurna', kita dapat berargumen bahwa jika bilangan itu
kuadrat, \( p=n^2 \), maka ia dapat difaktorkan sebagai
\( p=n\cdot n \) dengan \( n<p \), sehingga tidak prima
(\( p=0 \) atau \( p=1 \) tidak memenuhi \( n<p \), tetapi keduanya
bukan bilangan prima).

% Note two things about this statement form.

% First, an `if \( P \) then \( Q \)' result can sometimes be improved
% by weakening \( P \) or strengthening \( Q \).
% Thus, `if a number is divisible by \( p^2 \) then its square is also
% divisible by \( p^2 \)' could be upgraded either by relaxing its
% hypothesis:~`if a number is divisible by \( p \) then its square
% is divisible by \( p^2 \)', or by tightening its conclusion:~`if
% a number is divisible by \( p^2 \) then its square is divisible by
% \( p^4 \)'.

% Second,
% after showing `if \( P \) then \( Q \)' then a good next step is to look into
% whether there are cases where \( Q \) holds but \( P \) does not.
% The idea is to better understand the relationship between \( P \) and
% \( Q \) with an eye toward strengthening the proposition.



\startword{Pernyataan ekuivalen}
\index{pernyataan ekuivalen}\index{proposisi!ekuivalen}
Kadang-kadang bukan hanya \( P \) yang mengimplikasikan \( Q \), tetapi
\( Q \) juga mengimplikasikan \( P \).
Hubungan ini dapat dinyatakan dengan beberapa cara:
`\( P \) jika dan hanya jika \( Q \)', `\( P \) jhj \( Q \)',
`\( P \) dan \( Q \) ekuivalen secara logis',
`\( P \) perlu dan cukup untuk menghasilkan \( Q \)', atau
`\( P\iff Q \)'.
Sebagai contoh, `suatu bilangan bulat habis dibagi sepuluh jika dan hanya
jika angka terakhirnya adalah $0$'.

Walaupun dalam argumen sederhana rantai seperti
``\( P \) jika dan hanya jika $R$, yang berlaku jika dan hanya jika
$S$ \ldots'' mungkin praktis, untuk membuktikan bahwa dua pernyataan
ekuivalen kita lebih sering membuktikan kedua arah
`jika \( P \), maka \( Q \)' dan `jika \( Q \), maka \( P \)' secara terpisah.








\appendsection{Kuantor}\index{kuantor}%
Bandingkan pernyataan-pernyataan tentang bilangan asli berikut:
`terdapat bilangan asli \( x \) sehingga \( x \) habis dibagi \( x^2 \)'
benar, sedangkan
`untuk setiap bilangan asli \( x \), \( x \) habis dibagi \( x^2 \)'
salah.
Awalan `terdapat' dan `untuk setiap' disebut
\definend{kuantor}.\index{kuantor}

\startword{Untuk setiap}
Awalan `untuk setiap' adalah
\definend{kuantor universal},\index{kuantor!universal}
yang dilambangkan dengan \( \forall \).

Cara paling langsung untuk membuktikan bahwa suatu pernyataan berlaku
dalam semua kasus adalah membuktikannya untuk setiap kasus.
Jadi, untuk menunjukkan bahwa `setiap bilangan yang habis dibagi \( p \)
memiliki kuadrat yang habis dibagi \( p^2 \)', ambil satu bilangan berbentuk
\( pn \), lalu kuadratkan: \( (pn)^2=p^2n^2 \).
Ini adalah pembuktian dengan \definend{unsur tipikal}.
(Dalam argumen semacam ini, berhati-hatilah agar tidak mengasumsikan sifat
unsur tersebut selain sifat-sifat yang terdapat dalam hipotesis.
Argumen berikut salah:
``Jika \( n \) habis dibagi suatu bilangan prima, katakanlah \( 2 \),
sehingga \( n=2k \) untuk suatu bilangan asli~$k$, maka
\( n^2=(2k)^2=4k^2 \), sehingga kuadrat $n$ habis dibagi kuadrat bilangan
prima itu.''
Argumen tersebut hanya membuktikan kasus khusus \( p=2 \), bukan semua
nilai~\( p \).
Bandingkan dengan argumen yang benar:
``Jika \( n \) habis dibagi suatu bilangan prima sehingga \( n=pk \)
untuk suatu bilangan asli~$k$, maka \( n^2=(pk)^2=p^2k^2 \), sehingga
kuadrat $n$ habis dibagi kuadrat bilangan prima itu.'')




\startword{Terdapat}
Awalan `terdapat' adalah
\definend{kuantor eksistensial},\index{kuantor!eksistensial}
yang dilambangkan dengan \( \exists \).

%This quantifier is in some ways the opposite of `for all'.
%For instance, contrast these two definitions of primality
%of an integer \( p \): (i)~for all \( n \), if
%\( n \) is not \( 1 \) and \( n \) is not \( p \) then \( n \)
%does not divide \( p \), and
%(ii)~it is not the case that there exists an \( n \)
%(with \( n\neq 1 \) and \( n\neq p \)) such that \( n \) divides \( p \).

Kita dapat membuktikan proposisi keberadaan dengan memberikan sesuatu yang
memenuhi sifat tersebut. Sebagai contoh, untuk menyelesaikan pertanyaan
tentang keprimaan \( 2^{2^5}+1 \), Euler menunjukkan pembagi
\( 641 \)\cite{Sandifer}.
Namun, ada pula pembuktian yang menunjukkan bahwa sesuatu ada tanpa
menjelaskan cara menemukannya; argumen Euclid pada subbagian berikutnya
menunjukkan bahwa terdapat tak hingga banyak bilangan prima tanpa memberikan
rumus yang menamai bilangan-bilangan itu.
% In general, while demonstrating existence is better than nothing,
% giving an example is better, and an
% exhaustive list of all instances is ideal.

Akhirnya, setelah bertanya
``Apakah ada?''\ %\spacefactor=1000 
kita sering bertanya ``Ada berapa?''
Dengan kata lain, pertanyaan tentang ketunggalan sering muncul bersama
pertanyaan tentang keberadaan.
Kadang-kadang kedua argumen lebih sederhana jika dipisahkan. Karena itu,
perhatikan bahwa sebagaimana membuktikan keberadaan tidak membuktikan
ketunggalan, membuktikan ketunggalan pun tidak membuktikan keberadaan.










\appendsection{Teknik Pembuktian}
\index{teknik pembuktian|(}%
Ada banyak cara untuk membuktikan pernyataan matematika.
Di sini kita akan menguraikan dua teknik yang sering digunakan, tetapi
mungkin tidak langsung terasa alami bahkan bagi orang yang terbiasa
berpikir teknis.
 

\startword{Induksi}
\index{induksi matematika}
Banyak pembuktian bersifat iteratif,
``Inilah alasan pernyataan tersebut benar untuk bilangan \( 0 \),
lalu dari sana berlaku untuk \( 1 \), kemudian \( 2 \), dan seterusnya.''
Pembuktian semacam ini disebut pembuktian dengan
\definend{induksi matematika}.\index{induksi matematika}\index{induksi}
Kita akan melihat dua contoh.

Pertama, kita akan membuktikan bahwa \( 1+2+3+\dots+n=n(n+1)/2 \).
Rumus ini memiliki variabel bilangan asli bebas $n$. Artinya, menetapkan
$n$ sama dengan $1$, $2$, dan seterusnya menghasilkan suatu keluarga kasus:
kasus pertama $1=1(2)/2$, kasus kedua $1+2=2(3)/2$, dan seterusnya.
Pembuktian induksi kita melibatkan pernyataan dengan satu variabel bilangan
asli bebas.

Setiap pembuktian seperti itu mempunyai dua langkah.
Pada \definend{langkah dasar}\index{langkah dasar, pembuktian induksi}
kita menunjukkan bahwa pernyataan berlaku untuk suatu bilangan awal
$i\in \N$. Biasanya langkah ini merupakan verifikasi rutin.
Langkah kedua, yaitu
\definend{langkah induksi},\index{langkah induksi, pembuktian induksi}
lebih halus; kita akan menunjukkan bahwa implikasi berikut berlaku:
\begin{equation*}
  \begin{tabular}{l}
    \text{Jika pernyataan berlaku dari $n=i$ sampai dengan~$n=k$} 
     \\
    \text{\ maka pernyataan juga berlaku pada kasus $n=k+1$}
  \end{tabular}
  \tag{$*$}
\end{equation*}
(baris pertama adalah
\definend{hipotesis induksi}\index{induksi!hipotesis induksi}\index{hipotesis induksi}).
Menyelesaikan kedua langkah tersebut membuktikan bahwa pernyataan benar
untuk semua bilangan asli yang lebih besar daripada atau sama dengan $i$.

Untuk pernyataan tentang jumlah $n$ bilangan awal, intuisi di balik prinsip
ini adalah sebagai berikut. Pertama, langkah dasar memverifikasi secara
langsung pernyataan untuk bilangan awal $n=1$. Kemudian, karena langkah
induksi membuktikan implikasi~($*$) untuk setiap $k$, penerapan implikasi
tersebut pada $k=1$ menunjukkan bahwa pernyataan benar untuk $n=2$.
Setelah pernyataan terbukti untuk $1$ dan $2$, terapkan ($*$) lagi untuk
menyimpulkan bahwa pernyataan benar untuk $n=3$.
Dengan cara ini, kita melangkah terus hingga mencakup semua $n\geq 1$.

Berikut pembuktian \( 1+2+3+\dots+n=n(n+1)/2 \), dengan paragraf terpisah
untuk langkah dasar dan langkah induksi.
\begin{quote}\small
Untuk langkah dasar, kita menunjukkan bahwa rumus berlaku ketika \( n=1 \).
Ini mudah; jumlah satu bilangan asli pertama sama dengan \( 1(1+1)/2 \).

Untuk langkah induksi, andaikan hipotesis induksi bahwa rumus berlaku untuk
\( n=1 \), \( n=2 \), \ldots, \( n=k \) dengan $k\geq 1$.
Dengan kata lain, andaikan
$1=1(1)/2$, $1+2=2(3)/2$, $1+2+3=3(4)/2$, sampai
$1+2+\cdots+k=k(k+1)/2$.
Dengan demikian, rumus juga berlaku pada kasus \( n=k+1 \):
\begin{equation*}
  1+2+\cdots+k+(k+1)
  =
  \frac{k(k+1)}{2}+(k+1)
  =
  \frac{(k+1)(k+2)}{2}
\end{equation*}
(kesamaan pertama mengikuti hipotesis induksi).
\end{quote}
% The base case shows that the formula is true holds for \( n=1 \).
% The inductive step shows that,  
% because the formula holds for \( 1 \), it also holds for \( n=2 \);
% The inductive step also shows that,  
% because the formula holds for \( 1 \) and \( 2\), it holds for \( 3=2 \);
% Continuing in this way, we see that the statement holds
% for any natural number greater than or equal to \( 1 \).

Berikut contoh lain yang membuktikan bahwa setiap bilangan bulat yang lebih
besar daripada atau sama dengan \( 2 \) merupakan hasil kali bilangan prima.
\begin{quote}\small
Langkah dasarnya mudah: \( 2 \) adalah hasil kali satu bilangan prima.

Untuk langkah induksi, andaikan masing-masing dari \( 2,3,\ldots,k \)
merupakan hasil kali bilangan prima, dengan tujuan menunjukkan bahwa
\( k+1 \) juga merupakan hasil kali bilangan prima.
Ada dua kemungkinan.
Pertama, jika \( k+1 \) tidak habis dibagi bilangan yang lebih kecil daripada
dirinya, maka \( k+1 \) adalah bilangan prima dan dengan demikian merupakan
hasil kali bilangan prima.
Kemungkinan kedua adalah \( k+1 \) habis dibagi suatu bilangan yang lebih
kecil daripada dirinya. Berdasarkan hipotesis induksi, faktor-faktornya dapat
ditulis sebagai hasil kali bilangan prima.
Dalam kedua kasus, \( k+1 \) dapat ditulis sebagai hasil kali bilangan prima.
\end{quote}





\startword{Kontradiksi}
\index{kontradiksi}
Teknik pembuktian lain menunjukkan bahwa sesuatu benar dengan menunjukkan
bahwa sesuatu itu tidak mungkin salah.
Pembuktian dengan kontradiksi mengandaikan proposisi tersebut salah, lalu
menurunkan suatu pertentangan dengan fakta yang telah diketahui.

Contoh klasik pembuktian dengan kontradiksi adalah argumen Euclid bahwa
terdapat tak hingga banyak bilangan prima.
\begin{quote}\small
Andaikan hanya ada berhingga banyak bilangan prima \( p_1,\dots,p_k \).
Perhatikan bilangan \( p_1\cdot p_2\dots p_k +1 \).
Tidak satu pun bilangan prima dalam daftar yang dianggap lengkap itu membagi
bilangan tersebut tanpa sisa, sebab masing-masing menyisakan \( 1 \).
Namun, setiap bilangan merupakan hasil kali bilangan prima, sehingga keadaan
ini mustahil. Jadi, banyak bilangan prima tidak mungkin berhingga.
\end{quote}

Contoh lain adalah pembuktian bahwa \( \sqrt{2} \) bukan bilangan rasional.
\begin{quote}\small
Andaikan \( \sqrt{2}=m/n \), sehingga $2n^2=m^2$.
Keluarkan semua faktor \( 2 \), sehingga diperoleh
\( n=2^{k_n}\cdot \hat{n} \)
dan
\( m=2^{k_m}\cdot \hat{m} \).
Tuliskan kembali persamaan tersebut.
\begin{equation*}
  2\cdot (2^{k_n}\cdot \hat{n})^2
  =
  (2^{k_m}\cdot \hat{m})^2
\end{equation*}
Teorema Faktorisasi Prima menyatakan bahwa banyak faktor \( 2 \) pada kedua
ruas harus sama. Namun, ruas kiri memiliki banyak faktor yang ganjil,
\( 1+2k_n \), sedangkan ruas kanan memiliki banyak faktor yang genap,
\( 2k_m \). Ini suatu kontradiksi, sehingga tidak mungkin ada bilangan
rasional yang kuadratnya sama dengan \( 2 \).
\end{quote}

% Both of these examples aimed to prove something doesn't exist.
% A negative proposition often suggests a proof by contradiction.
\index{teknik pembuktian|)}













\appendsection{Himpunan, Fungsi, dan Relasi}
Materi berikut menjadi latar belakang dan kosakata bagi seluruh pembahasan
yang kita kembangkan.

\startword{Himpunan}
\index{himpunan}
Matematikawan sering bekerja dengan koleksi.
Jenis koleksi yang paling umum digunakan adalah \definend{himpunan}.\index{himpunan}
Himpunan dicirikan oleh Prinsip Ekstensionalitas: dua himpunan yang memiliki
unsur-unsur sama adalah sama.
Karena itu, urutan unsur tidak penting,
\( \set{2,\pi}=\set{\pi,2} \), dan pengulangan dilebur,
\( \set{7,7}=\set{7} \).

Kita dapat mendeskripsikan himpunan dengan daftar di antara kurung kurawal,
\( \set{1,4,9,16} \), seperti dalam paragraf sebelumnya, atau dengan
notasi pembentuk himpunan
\( \set{x\suchthat x^5-3x^3+2=0} \), yang dibaca ``himpunan semua \( x \)
sedemikian sehingga \ldots''.
Kita menamai himpunan dengan huruf Romawi kapital; misalnya, himpunan
bilangan prima adalah \( P=\set{2,3,5,7,11,\ldots\,} \).
Beberapa himpunan begitu penting sehingga namanya telah ditetapkan, seperti
bilangan riil \( \Re \) dan bilangan kompleks \( \C \).
Untuk menyatakan bahwa sesuatu merupakan
\definend{unsur},\index{unsur}\index{himpunan!unsur}
atau \definend{anggota}\index{anggota}\index{himpunan!anggota} suatu himpunan,
kita memakai `\(\in\)'. Jadi \( 7\in\set{3,5,7} \), sedangkan
\( 8\not\in\set{3,5,7} \).

Kita mengatakan bahwa \( A \) adalah \definend{subhimpunan} dari \( B \),
ditulis $A\subseteq B$, apabila $x\in A$ mengimplikasikan $x\in B$.
Dalam buku ini, kita memakai `\( \subset \)' untuk relasi
\definend{subhimpunan sejati}\index{himpunan!subhimpunan sejati}\index{subhimpunan sejati}\index{himpunan!subhimpunan}\index{subhimpunan!sejati} %
bahwa \( A \) adalah subhimpunan \( B \), tetapi \( A\neq B \)
(sebagian penulis memakai lambang ini untuk semua jenis subhimpunan,
baik sejati maupun tidak).
Sebagai contoh, \( \set{2,\pi}\subset\set{2,\pi,7} \).
Lambang-lambang ini dapat dibalik; misalnya,
\( \set{2,\pi,5}\supset\set{2,5} \).

Berdasarkan Ekstensionalitas, untuk membuktikan dua himpunan sama,
\( A=B \), tunjukkan bahwa keduanya memiliki anggota yang sama.
Sering kali kita melakukannya dengan menunjukkan inklusi timbal balik,
\index{inklusi timbal balik}\index{himpunan!inklusi timbal balik}
yaitu \( A\subseteq B \) sekaligus \( A\supseteq B \).
Pembuktian seperti itu mempunyai satu bagian yang menunjukkan bahwa jika
$x\in A$, maka $x\in B$, dan bagian kedua yang menunjukkan bahwa jika
$x\in B$, maka $x\in A$.

Himpunan tanpa anggota disebut
\definend{himpunan kosong}\index{himpunan kosong}\index{himpunan!kosong}
\( \set{} \), yang dilambangkan dengan \( \emptyset \).
Setiap himpunan memiliki himpunan kosong sebagai subhimpunan berdasarkan
sifat `benar secara hampa' dalam definisi implikasi.




\startword{Diagram}
Kita menggambarkan operasi dasar himpunan dengan
\definend{diagram Venn}.\index{diagram Venn}
Diagram berikut menunjukkan \( x\in P \).
\begin{center}
  \includegraphics{appen/mp/appen.8}
\end{center}
Persegi panjang terluar memuat semesta $\Omega$ dari semua objek yang sedang
dibahas. Misalnya, dalam pernyataan tentang bilangan riil, persegi panjang
tersebut memuat semua anggota $\R$.
Himpunan digambarkan sebagai lingkaran yang mengurung anggota-anggotanya.

Berikut diagram untuk \( P\subseteq Q \).
Diagram ini menunjukkan bahwa jika \( x\in P \), maka \( x\in Q \).
\begin{center}
  \includegraphics{appen/mp/appen.4}
\end{center}




\startword{Operasi Himpunan}
\definend{Gabungan}\index{gabungan himpunan}\index{himpunan!gabungan}
dua himpunan adalah
\( P\union Q=\set{x\suchthat \text{$(x\in P)$ atau $(x\in Q)$}} \).
Diagram menunjukkan bahwa suatu unsur berada dalam gabungan jika unsur itu
berada dalam sedikitnya salah satu himpunan.
\begin{center}
  \includegraphics{appen/mp/appen.3}
\end{center}
\definend{Irisan}\index{irisan himpunan}\index{himpunan!irisan} adalah
\( P\intersection Q=\set{x\suchthat \text{$(x\in P)$ dan $(x\in Q)$}} \).
\begin{center}
  \includegraphics{appen/mp/appen.2}
\end{center}
\definend{Komplemen}\index{komplemen}\index{himpunan!komplemen}
himpunan~\( P \) adalah
\( P^{\text{comp}}=\set{x\in\Omega\suchthat x\not\in P} \)
\begin{center}
  \includegraphics{appen/mp/appen.1}
\end{center}




\startword{Multihimpunan}\index{multihimpunan}
Seperti dijelaskan di atas, himpunan adalah koleksi yang urutannya tidak
penting, sehingga $\set{2,\pi}$ dan $\set{\pi,2}$ sama, dan yang
pengulangannya dilebur, sehingga $\set{7,7}$ dan $\set{7}$ sama.

Koleksi yang menyerupai himpunan karena urutannya tidak penting, tetapi
pengulangannya tidak dilebur, disebut \definend{multihimpunan}.
(Perhatikan bahwa kita memakai notasi kurung kurawal $\set{\ldots}$ yang
sama seperti untuk himpunan.)
Jadi multihimpunan $\set{1,2,2}$ berbeda dari multihimpunan $\set{1,2}$.
Karena urutan tidak penting, multihimpunan $\set{1,2,2}$,
$\set{2,1,2}$, dan~$\set{2,2,1}$ semuanya sama.
Dalam buku ini multihimpunan hanya disebutkan dalam sebuah catatan, sehingga
pembahasan tentang subhimpunan, gabungan, atau irisannya berada di luar
cakupan kita.



\startword{Barisan}
\index{barisan}
Selain himpunan dan multihimpunan, kita juga memakai koleksi yang urutannya
penting dan pengulangannya tidak dilebur.
Koleksi ini disebut \definend{barisan},\index{barisan} dan dilambangkan dengan
kurung sudut: \( \sequence{2,3,7}\neq\sequence{2,7,3} \).
Barisan dengan panjang \( 2 \) disebut
\definend{pasangan terurut},\index{pasangan terurut}\index{pasangan!terurut}
dan sering ditulis dengan tanda kurung: \( (\pi,3) \).
Kadang-kadang kita juga mengatakan `tripel terurut', `\( 4 \)-tupel
terurut', dan seterusnya.
Himpunan semua \( n \)-tupel terurut dari unsur-unsur himpunan \( A \)
dilambangkan dengan \( A^n \).
Jadi \( \Re^2 \) adalah himpunan pasangan bilangan riil.




\startword{Fungsi}
\index{fungsi}
\definend{Fungsi}\index{fungsi}
atau \definend{pemetaan}\index{pemetaan} $\map{f}{D}{C}$ adalah
pengaitan antara masukan
\definend{argumen}\index{argumen, fungsi}\index{fungsi!argumen}
$x\in D$ dan keluaran
\definend{nilai}\index{nilai, fungsi}\index{fungsi!nilai}
$f(x)\in C$, dengan syarat bahwa fungsi itu harus
\definend{terdefinisi dengan baik},\index{fungsi!terdefinisi dengan baik}%
\index{terdefinisi dengan baik}
yaitu \( x \) cukup untuk menentukan \( f(x) \).
Dengan kata lain, syaratnya adalah bahwa jika \( x_1=x_2 \), maka
\( f(x_1)=f(x_2) \).

Himpunan semua argumen~$D$ disebut
\definend{domain}\index{domain}\index{fungsi!domain} dari \( f \),
sedangkan himpunan nilai keluarannya disebut
\definend{daerah hasil} $\rangespace{f}$.\index{daerah hasil}\index{fungsi!daerah hasil}
Sering kali kita tidak bekerja dengan daerah hasil, melainkan dengan suatu
superhimpunan yang lebih nyaman, yaitu
\definend{kodomain}~$C$.\index{fungsi!kodomain}\index{kodomain}
Misalnya, kita dapat mendeskripsikan fungsi kuadrat sebagai
$\map{s}{\R}{\R}$ alih-alih $\map{s}{\R}{\R^+\union \set{0}}$.

Kita menggambarkan fungsi dengan
\definend{diagram kacang}.\index{diagram kacang}\index{fungsi!diagram kacang}
\begin{center}
  \includegraphics{appen/mp/appen.9}
\end{center}
Gumpalan di sebelah kiri adalah domain, sedangkan gumpalan di sebelah kanan
adalah kodomain.
Fungsi tersebut mengaitkan tiga titik dalam domain dengan tiga titik dalam
kodomain.
Perhatikan bahwa menurut definisi fungsi, setiap titik dalam domain dikaitkan
dengan tepat satu titik dalam kodomain, tetapi kebalikannya tidak harus benar.

Pengaitannya dapat bersifat sebarang; tidak diperlukan rumus atau algoritma,
walaupun dalam buku ini biasanya ada salah satunya.
Kita sering memakai $y$ untuk melambangkan $f(x)$.
Kita juga memakai notasi \( x\mapsunder{f}16x^2-100 \), yang dibaca
`\( x \) dipetakan oleh \( f \) ke \( 16x^2-100 \)' atau
`\( 16x^2-100 \) adalah
\definend{citra}\index{citra, oleh fungsi}\index{fungsi!citra}
dari \( x \)'.

Pemetaan seperti \( x\mapsto \sin(1/x) \) merupakan gabungan pemetaan yang
lebih sederhana, yakni \( g(y)=\sin(y) \) diterapkan pada citra
\( f(x)=1/x \).
\definend{Komposisi}\index{komposisi}\index{fungsi!komposisi}
\( \map{g}{Y}{Z} \) dengan \( \map{f}{X}{Y} \) adalah pemetaan yang
mengirim \( x\in X \) ke \( g(\,f(x)\,)\in Z \).
Komposisi ini dilambangkan dengan \( \map{\composed{g}{f}}{X}{Z} \).
Definisi ini hanya bermakna apabila daerah hasil \( f \) merupakan
subhimpunan dari domain \( g \).

\definend{Pemetaan identitas}\index{identitas!fungsi}\index{fungsi!identitas}
\( \map{\identity}{Y}{Y} \), yang didefinisikan oleh
\( \identity(y)=y \), memiliki sifat bahwa untuk setiap
\( \map{f}{X}{Y} \), komposisi \( \composed{\identity}{f} \) sama dengan
\( f \).
Jadi, terhadap komposisi fungsi, pemetaan identitas memainkan peran yang sama
seperti bilangan \( 0 \) dalam penjumlahan bilangan riil atau \( 1 \) dalam
perkalian.

Sejalan dengan analogi tersebut, kita mendefinisikan
\definend{invers kiri}\index{invers!fungsi!kiri}\index{invers!kiri}\index{invers kiri}
dari pemetaan \( \map{f}{X}{Y} \) sebagai fungsi
\( \map{g}{\text{daerah hasil}(f)}{X} \) sedemikian sehingga
\( \composed{g}{f} \) adalah pemetaan identitas pada \( X \).
\definend{Invers kanan}\index{invers!fungsi!kanan}\index{invers!kanan}\index{invers kanan}
dari \( f \) adalah pemetaan \( \map{h}{Y}{X} \) sedemikian sehingga
\( \composed{f}{h} \) adalah identitas.

Untuk beberapa fungsi $f$, terdapat pemetaan yang sekaligus merupakan invers
kiri dan invers kanan dari \( f \).
Jika pemetaan seperti itu ada, maka pemetaan itu tunggal, sebab jika
\( g_1 \) dan \( g_2 \) sama-sama memiliki sifat tersebut, maka
\( g_1(x)=\composed{g_1}{ (\composed{f}{g_2}) }\,(x)
         =\composed{ (\composed{g_1}{f}) }{g_2}\,(x)
         =g_2(x) \)
(kesamaan di tengah berasal dari sifat asosiatif komposisi fungsi).
Karena itu, kita menyebutnya \definend{invers dua sisi}, atau cukup
\definend{``invers''},\index{invers}\index{invers!dua sisi}\index{fungsi!invers}\index{invers!fungsi}\index{fungsi invers}\index{penginversan}
dan melambangkannya dengan \( f^{-1} \).
Sebagai contoh, invers fungsi \( \map{f}{\Re}{\Re} \) yang diberikan oleh
\( f(x)=2x-3 \) adalah fungsi \( \map{f^{-1}}{\Re}{\Re} \) yang diberikan
oleh \( f^{-1}(x)=(x+3)/2 \).

Notasi superskrip untuk invers fungsi, `\( f^{-1} \)', merupakan bagian dari
suatu pola yang lebih luas.
Fungsi yang domain dan kodomainnya sama, $\map{f}{X}{X}$, dapat diiterasikan,
sehingga kita dapat meninjau komposisi $f$ dengan dirinya sendiri:
\( \composed{f}{f} \), \( \composed{f}{\composed{f}{f}} \), dan seterusnya.
Kita menulis $\composed{f}{f}$ sebagai \( f^2 \) dan
$\composed{\composed{f}{f}}{f}$ sebagai \( f^3 \), dan seterusnya.
Perhatikan bahwa aturan pangkat yang dikenal untuk bilangan riil berlaku:
\( \composed{f^i}{f^j}=f^{i+j} \) dan \( (f^i)^j=f^{i\cdot j} \).
Jika \( f \) dapat diinverskan, penulisan \( f^{-1} \) untuk inversnya dan
\( f^{-2} \) untuk invers dari \( f^2 \), dan seterusnya, membuat aturan
pangkat tersebut tetap berlaku karena kita mendefinisikan \( f^0 \) sebagai
pemetaan identitas.

Definisi fungsi mensyaratkan bahwa untuk setiap masukan terdapat tepat satu
nilai keluaran yang terkait.
Jika fungsi $\map{f}{D}{C}$ memiliki sifat tambahan bahwa untuk setiap nilai
keluaran terdapat sedikitnya satu nilai masukan yang terkait\Dash dengan kata
lain, kodomain $f$ sama dengan daerah hasilnya,
\( C=\rangespace{f} \)\Dash maka fungsi itu
\definend{surjektif}.\index{fungsi!surjektif}\index{fungsi surjektif}
\begin{center}
  \includegraphics{appen/mp/appen.14}
\end{center}
Suatu fungsi memiliki invers kanan jika dan hanya jika fungsi itu surjektif.
(Fungsi $f$ pada gambar di atas mempunyai invers kanan $\map{g}{C}{D}$ yang
diperoleh dengan mengikuti panah ke arah sebaliknya, dari kanan ke kiri.
Untuk titik kodomain teratas, pilih salah satu dari dua panah yang dapat
diikuti. Dengan demikian, menerapkan $g$ lalu $f$ membawa setiap unsur
$y\in C$ kembali kepada dirinya sendiri, sehingga komposisinya adalah fungsi
identitas.)

Jika fungsi $\map{f}{D}{C}$ memiliki sifat bahwa untuk setiap nilai keluaran
terdapat paling banyak satu nilai masukan yang terkait\Dash yakni tidak ada
dua argumen yang memiliki citra sama, sehingga \( f(x_1)=f(x_2) \)
mengimplikasikan \( x_1=x_2 \)\Dash maka fungsi tersebut
\definend{injektif}.\index{fungsi!injektif}\index{fungsi injektif}
Diagram kacang sebelumnya menggambarkan sifat ini.
\begin{center}
  \includegraphics{appen/mp/appen.9}
\end{center}
Suatu fungsi memiliki invers kiri jika dan hanya jika fungsi itu injektif.
(Pada gambar, definisikan $\map{g}{C}{D}$ dengan mengikuti panah ke arah
sebaliknya bagi $y\in C$ yang berada di ujung suatu panah, dan mengirim titik
tersebut ke unsur sebarang dalam $D$ pada kasus lainnya.
Menerapkan $f$ lalu~$g$ pada unsur-unsur $D$ kemudian bertindak sebagai
identitas.)

Berdasarkan paragraf-paragraf sebelumnya, suatu pemetaan memiliki invers dua
sisi jika dan hanya jika pemetaan itu surjektif sekaligus injektif.
Fungsi seperti itu merupakan suatu
\definend{korespondensi}.\index{korespondensi}\index{fungsi!korespondensi}
Fungsi tersebut mengaitkan tepat satu unsur domain dengan setiap unsur
kodomain.
Karena komposisi pemetaan injektif bersifat injektif dan komposisi pemetaan
surjektif bersifat surjektif, komposisi korespondensi juga merupakan
korespondensi.

Kadang-kadang kita ingin memperkecil domain suatu fungsi.
Sebagai contoh, kita dapat mengambil fungsi \( \map{f}{\Re}{\Re} \) yang
diberikan oleh \( f(x)=x^2 \), lalu, agar memiliki invers, membatasi argumen
masukan pada bilangan riil tak negatif,
\( \map{\hat{f}}{\Re^+\union\set{0}}{\Re} \).
Fungsi \( \hat{f} \) kemudian disebut
\definend{restriksi}\index{fungsi!restriksi}\index{restriksi} dari \( f \)
pada domain yang lebih kecil.








\startword{Relasi}
\index{relasi}
Beberapa objek matematika yang dikenal, seperti `\( < \)' atau `\( = \)',
paling alami dipahami sebagai relasi antarobjek.
\definend{Relasi biner} pada himpunan \( A \) adalah himpunan pasangan
terurut dari unsur-unsur \( A \).
Sebagai contoh, beberapa unsur himpunan yang merupakan relasi `$<$' pada
bilangan bulat adalah \( (3,5) \), \( (3,7) \), dan \( (1,100) \).
Relasi biner lain pada bilangan bulat adalah kesamaan; relasi ini merupakan
himpunan \( \set{\ldots,(-1,-1),(0,0),(1,1),\ldots} \).
Contoh lain adalah `berjarak kurang dari \( 10 \)', yaitu himpunan
\( \set{(x,y)\suchthat |x-y|<10} \).
Beberapa anggota relasi ini adalah \( (1,10) \), \( (10,1) \), dan
\( (42,44) \). Baik \( (11,1) \) maupun \( (1,11) \) bukan anggotanya.

Contoh-contoh tersebut memperlihatkan luasnya cakupan definisi ini.
Segala jenis hubungan tercakup, misalnya `kedua bilangan genap' atau
`bilangan pertama diperoleh dengan membalik urutan angka bilangan kedua'.




\startword{Relasi Ekuivalensi}
\index{relasi!ekuivalensi}\index{relasi ekuivalensi}\index{ekuivalensi}
Kita perlu menyatakan bahwa dua objek serupa dalam suatu hal.
Keduanya tidak identik, tetapi saling berelasi
(misalnya, dua bilangan bulat yang memberikan sisa sama ketika dibagi \( 2 \)).

Relasi biner \( \set{(a,b),\ldots} \) disebut
\definend{relasi ekuivalensi}\index{ekuivalensi!relasi}\index{relasi!ekuivalensi}
apabila memenuhi
(1)~\definend{refleksivitas}:\index{refleksivitas, relasi}\index{relasi!refleksif}
     setiap objek berelasi dengan dirinya sendiri,
(2)~\definend{simetri}:\index{simetri, relasi}\index{relasi!simetris}
     jika \( a \) berelasi dengan \( b \), maka \( b \) berelasi dengan
     \( a \), dan
(3)~\definend{transitivitas}:\index{transitivitas, relasi}\index{relasi!transitif}
     jika \( a \) berelasi dengan \( b \) dan \( b \) berelasi dengan
     \( c \), maka \( a \) berelasi dengan \( c \).
Beberapa contoh pada bilangan bulat: `\( = \)' merupakan relasi ekuivalensi;
`\( < \)' tidak memenuhi simetri; `bertanda sama' merupakan relasi
ekuivalensi; sedangkan `berjarak kurang dari \( 10 \)' tidak memenuhi
transitivitas.






\startword{Partisi}
\index{partisi|(}
Dalam relasi `bertanda sama'
\( \set{(1,3),(-5,-7),(0,0),\ldots} \), terdapat tiga jenis pasangan:
pasangan yang kedua bilangannya positif, pasangan yang keduanya negatif,
dan satu pasangan yang keduanya nol.
Jadi, setiap bilangan bulat berada tepat dalam salah satu dari tiga kelas:
positif, negatif, atau nol.

\definend{Partisi}\index{partisi} dari himpunan \( \Omega \) adalah koleksi
subhimpunan \( \set{S_0,S_1,S_2,\ldots} \) sedemikian sehingga setiap unsur
\( \Omega \) menjadi anggota suatu bagian,
\( S_0\union S_1\union S_2\union\cdots{}=\Omega \), dan bagian-bagian yang
beririsan adalah sama: jika \( S_i\intersection S_j\neq\emptyset \), maka
$S_i=S_j$.
Bayangkan \( \Omega \) diuraikan menjadi bagian-bagian yang tidak tumpang tindih.
\begin{center}
  \includegraphics{appen/mp/appen.10}
\end{center}
Dengan demikian, paragraf sebelumnya menyatakan bahwa `bertanda sama'
memartisi bilangan bulat menjadi himpunan bilangan positif, himpunan bilangan
negatif, dan himpunan yang hanya berisi nol.
Demikian pula, relasi ekuivalensi `=' memartisi bilangan bulat menjadi
himpunan-himpunan beranggota satu.

Contoh lain adalah himpunan untai yang terdiri atas suatu bilangan, diikuti
garis miring, lalu bilangan tak nol,
$\Omega=\set{n/d\suchthat \text{$n,d\in\Z$ dan $d\neq 0$}}$.
Definisikan $S_{n,d}$ dengan: $\hat{n}/\hat{d}\in S_{n,d}$ jika
$\hat{n}d=n\hat{d}$.
Mengecek bahwa ini merupakan partisi dari $\Omega$ adalah pekerjaan rutin
(misalnya, perhatikan bahwa $S_{4,3}=S_{8,6}$).
Gambar berikut menunjukkan beberapa bagian, dengan mencantumkan dua dari tak
hingga banyak anggotanya pada setiap bagian.
\begin{center}
  \includegraphics{appen/mp/appen.11}
\end{center}

Setiap relasi ekuivalensi menginduksi suatu partisi, dan setiap partisi
diinduksi oleh suatu relasi ekuivalensi.
(Hal ini mudah diperiksa.)
Berikut dua contoh.

Tinjau relasi ekuivalensi pada dua bilangan bulat `memberikan sisa yang sama
ketika dibagi \( 2 \)', yaitu himpunan
\( P=\set{(-1,3),(2,4),(0,0),\ldots} \).
Dalam himpunan $P$ terdapat dua jenis pasangan: pasangan yang kedua anggotanya
genap dan pasangan yang kedua anggotanya ganjil.
Ekuivalensi ini menginduksi suatu partisi dengan bagian-bagian yang ditentukan
sebagai berikut: untuk setiap \( x \), definisikan himpunan bilangan yang
berelasi dengannya, \( S_x=\set{y\suchthat (x,y)\in P} \).
Bagian-bagiannya adalah
\( \set{\ldots,-3,-1,1,3,\ldots} \) dan
\( \set{\ldots,-2,0,2,4,\ldots} \).
Setiap bagian dapat dinamai dengan banyak cara; misalnya,
\( \set{\ldots,-3,-1,1,3,\ldots} \) adalah \( S_1 \) sekaligus \( S_{-3} \).

Sekarang tinjau partisi bilangan asli dengan dua bilangan berada dalam bagian
yang sama apabila keduanya menyisakan sisa yang sama ketika dibagi $10$,
yaitu apabila keduanya memiliki angka satuan yang sama.
Partisi ini diinduksi oleh relasi ekuivalensi $R$ yang didefinisikan sebagai
berikut: dua bilangan $n$ dan $m$ berelasi apabila keduanya berada dalam
bagian yang sama.
Sebagai contoh, $3$ berelasi dengan $33$, tetapi $3$ tidak berelasi dengan
$102$.
Memverifikasi ketiga syarat dalam definisi ekuivalensi adalah pekerjaan mudah.

Setiap bagian dari suatu partisi disebut \definend{kelas ekuivalensi}.%
\index{ekuivalensi!kelas}\index{kelas!ekuivalensi}
Kadang-kadang kita memilih satu unsur dari setiap kelas ekuivalensi sebagai
\definend{wakil kelas}.%
\index{ekuivalensi!wakil}\index{kelas!wakil}\index{wakil!kelas}
\begin{center}
  \includegraphics{appen/mp/appen.13}
\end{center}
Biasanya, ketika memilih wakil kita mempertimbangkan suatu aturan yang alami.
Dalam hal itu, wakil tersebut disebut wakil
\definend{kanonik}.%
\index{wakil alami}\index{wakil kanonik}%
\index{ekuivalensi!kelas!wakil kanonik}%
\index{wakil!kanonik}
Sebagai contoh, pecahan \( 3/5 \) dan \( 9/15 \) ekuivalen.
Dalam penggunaan sehari-hari, kita sering memilih pecahan dalam `bentuk
paling sederhana' atau `bentuk tersederhanakan', $3/5$, sebagai wakil kelas.
\begin{center}
  \includegraphics{appen/mp/appen.12}
\end{center}
\index{partisi|)}
%
%\end{document}
